\chapter{La bicapa: información par y transporte geométrico impar}
\label{chap:paridad-bicapa-barbero-rev6}
\index{entropía!bicapa}
\index{Barbero--Immirzi!paridad}
\index{hoja!intercambio}

Barbero--Immirzi y la entropía de las dos hojas no son dos cantidades
yuxtapuestas. Nacen de los mismos canales \(q_\pm\), pero conservan datos
distintos. La normalización probabilística olvida la escala común y mide la
distribución entre orientaciones; el funcional hexada--octada conserva el
signo de esa orientación. La primera lectura es par y la segunda impar.

\section{Distribución normalizada entre hojas}

Escribamos
\[
 x=A_{\rm rad},
 \qquad
 y=C^*_{\rm rad},
 \qquad
 q_-=e^{-x+y},
 \qquad
 q_+=e^{-x-y}.
\]
La suma \(q_-+q_+=2e^{-x}\cosh y\) permite definir
\[
 p_-=\frac{q_-}{q_-+q_+}
     =\frac{e^y}{2\cosh y},
 \qquad
 p_+=\frac{q_+}{q_-+q_+}
     =\frac{e^{-y}}{2\cosh y}.
\]
El factor común \(e^{-x}\) desaparece. La distribución conserva la asimetría
entre hojas, pero no la escala angular común.

\begin{definition}[Entropía bicapa]
\label{def:entropia-bicapa-rev6}
La entropía de la distribución orientada es
\[
 S_{\rm bi}(y)
 :=
 -p_-\log p_--p_+\log p_+
 =
 \log(2\cosh y)-y\tanh y.
\]
\end{definition}

\begin{proposition}[Paridad y máximo de la entropía bicapa]
\label{prop:paridad-entropia-bicapa-rev6}
La función \(S_{\rm bi}\) es par y satisface
\[
 S_{\rm bi}(-y)=S_{\rm bi}(y),
 \qquad
 S_{\rm bi}'(y)=-y\,\operatorname{sech}^2y,
 \qquad
 0<S_{\rm bi}(y)\leq\log2.
\]
El máximo \(\log2\) se alcanza únicamente en \(y=0\), cuando ambas hojas
tienen el mismo peso.
\end{proposition}

\begin{proof}
La paridad se sigue de que \(\cosh\) es par y \(y\tanh y\) también lo es.
La derivación directa produce
\[
 \frac{\dd}{\dd y}\log(2\cosh y)=\tanh y,
 \qquad
 \frac{\dd}{\dd y}(y\tanh y)
 =\tanh y+y\operatorname{sech}^2y.
\]
Por tanto \(S_{\rm bi}'(y)=-y\operatorname{sech}^2y\): la función crece en
\((-\infty,0)\), decrece en \((0,\infty)\) y toma en el origen el valor
\(\log2\). La positividad es la de la entropía de una distribución con dos
pesos estrictamente positivos.
\end{proof}

\section{Paridad del transporte hexada--octada}

Sea
\[
 s_n(q_-,q_+)=q_-^n-q_+^n.
\]
El intercambio de hojas es la involución
\[
 \mathsf J_{\rm sh}:(x,y)\longmapsto(x,-y),
\]
que permuta \(q_-\leftrightarrow q_+\). De aquí,
\[
 s_n\circ\mathsf J_{\rm sh}=-s_n.
\]
El funcional construido en el \cref{mdg:chap:barbero} es
\[
 \Gamma_{6\to8}(A,C^*)
 =
 \frac{\pi A C^*}{180}
{}+12\bigl(S_{90}(q_-)-S_{90}(q_+)\bigr)
-\bigl(S_{120}(q_-)-S_{120}(q_+)\bigr).
\]

\begin{theorem}[Teorema de paridad bicapa]
\label{thm:paridad-bicapa-barbero-rev6}
Bajo el intercambio \(C^*\mapsto-C^*\),
\[
 \boxed{S_{\rm bi}(-C^*_{\rm rad})
        =S_{\rm bi}(C^*_{\rm rad}),}
\]
\[
 \boxed{\Gamma_{6\to8}(A,-C^*)
        =-\Gamma_{6\to8}(A,C^*).}
\]
La misma bicapa produce, por tanto, una lectura informacional par y un
transporte geométrico impar.
\end{theorem}

\begin{proof}
La primera igualdad es la
\cref{prop:paridad-entropia-bicapa-rev6}. En la segunda, el término
\(\pi AC^*/180\) cambia de signo y cada diferencia de Lambert cambia de
signo porque se intercambian sus dos argumentos. La combinación completa es
impar.
\end{proof}

\section{Interpretación del operador constitutivo como acoplamiento entre vacío, campo y respuesta}

La entropía bicapa y \(\Gamma_{6\to8}\) no son dos nombres para el mismo
número. La primera cuantifica cuánto peso comparten las dos orientaciones
después de cancelar la escala común. El segundo conserva la orientación y la
incidencia del paso \(6\to8\). En el punto HMT,
\[
 S_{\rm bi}
 =0.69246069960358548\ldots\ {\rm nats},
 \qquad
 \Gamma_{6\to8}
 =0.27400767269445927\ldots.
\]
El vínculo entre información y geometría consiste en que ambas lecturas
factorizan por los mismos canales \(q_\pm\), no en igualar sus valores.

La banda partícula--antipartícula del
\cref{chap:banda-particula-virtual} reproduce la misma distinción de
paridad: la masa y la cantidad de correlación descienden al cociente de
banda, mientras carga, orientación y transporte cambian de signo. En la
acción de Cartan--Holst, \(\Gamma_{6\to8}\) ocupa la coordenada orientada;
en la lectura informacional, \(S_{\rm bi}\) cuenta la distribución entre
hojas. Así se explica por qué una sola estructura puede conservar
información e invertir orientación sin confundir ambas operaciones.
